% !TeX program = xelatex
% !BIB program = biber
% 中文简历；信息核对日期：2026-10-16。
% 来源：index/research/publications/teaching/awards/students/activities.jemdoc。
% 项目与早期任职细节保留自原简历；学院中文名称参照湖南大学教师介绍。
\documentclass[11pt,a4paper,roman]{moderncv}
\moderncvstyle{banking}
\moderncvcolor{purple}
\usepackage{xeCJK}
\setmainfont{texgyrepagella-regular.otf}[BoldFont=texgyrepagella-bold.otf,
  ItalicFont=texgyrepagella-italic.otf,BoldItalicFont=texgyrepagella-bolditalic.otf]
\setsansfont{texgyreheros-regular.otf}[BoldFont=texgyreheros-bold.otf,
  ItalicFont=texgyreheros-italic.otf,BoldItalicFont=texgyreheros-bolditalic.otf]
% 使用 TeX Live 自带字体，避免依赖系统中文字体。
\setCJKmainfont{FandolSong-Regular.otf}[BoldFont=FandolSong-Bold.otf,ItalicFont=FandolKai-Regular.otf]
\setCJKsansfont{FandolHei-Regular.otf}[BoldFont=FandolHei-Bold.otf]
\setCJKmonofont{FandolFang-Regular.otf}
\usepackage[left=22mm,right=22mm,top=20mm,bottom=20mm]{geometry}
\usepackage{needspace}
\usepackage{tabularx}
% 银行式模板原来的 l/r 列无法换行；长项目名称与编号需自动折行。
\renewcommand*{\cventry}[7][0.9em]{%
  \begin{tabularx}{\textwidth}{@{}X r@{}}%
    {\bfseries #4} & {\bfseries #5}\\%
    {\itshape #3\ifthenelse{\equal{#6}{}}{}{, #6}} & {\itshape #2}\\%
  \end{tabularx}%
  \ifx&#7&\else{\\\vbox{\small#7}}\fi%
  \par\addvspace{#1}}
\renewcommand*{\subsectionfont}{\large\bfseries\upshape}
\usepackage[backend=biber,style=numeric,sorting=none,maxbibnames=99,giveninits=false,
  defernumbers=true,isbn=false,url=true,doi=true]{biblatex}
\addbibresource{publications.bib}
% moderncv 在导言末尾加载 hyperref，随后再设置链接和中文元数据。
\AtEndPreamble{\hypersetup{colorlinks=true,linkcolor=color1,urlcolor=color1,
  pdftitle={魏恒峰中文简历},pdfauthor={魏恒峰},pdfsubject={学术简历；更新于2026年10月16日}}}
\urlstyle{same}
\setlength{\emergencystretch}{2em}
\renewcommand*{\mkbibnamegiven}[1]{\ifitemannotation{highlight}{\textbf{#1}}{#1}}
\renewcommand*{\mkbibnamefamily}[1]{\ifitemannotation{highlight}{\textbf{#1}}{#1}}
\renewcommand*{\bibfont}{\small}
\setlength{\bibitemsep}{0.65\baselineskip}
\AtEveryBibitem{\needspace{5\baselineskip}}
\name{魏恒峰}{}
\title{\Large 湖南大学 副教授}
\address{\href{https://csee.hnu.edu.cn/people/weihengfeng}{信息科学与工程学院}}{}{}
\phone[mobile]{13905194610}
\email{hfwei@hnu.edu.cn}
\homepage{hengxin.github.io}
\extrainfo{\href{https://scholar.google.com/citations?user=HK9-98EAAAAJ}{Google Scholar}
  \enspace | \enspace\href{https://github.com/hengxin}{GitHub}}

\begin{document}
\makecvtitle
\vspace{-1.5em}
\begin{center}\small 更新日期：2026 年 10 月 16 日\end{center}

\section{个人简介与研究方向}
2025 年 8 月起任湖南大学信息科学与工程学院副教授。
2009 年、2016 年于南京大学获得计算机科学与技术专业学士、博士学位；
2017 年 1 月至 2025 年 7 月先后在南京大学计算机学院、软件学院工作。

\medskip
研究领域包括\textbf{数据库系统、分布式计算与形式化方法}，
重点关注数据库事务一致性（隔离级别）的建模、规约与验证，
以及分布式数据一致性、复制数据类型和协议正确性。
相关研究使用 TLA+/TLAPS、Coq、SAT/SMT 等形式化工具与技术。
在 VLDB、SIGMOD、ICDE、SOSP、OOPSLA、ISSTA、ASE、PODC、USENIX ATC、TC、TPDS 等会议与期刊发表论文。

\section{教育经历}
\cventry{2009 年 9 月 -- 2016 年 12 月}{计算机科学与技术，博士}{南京大学}{计算机学院}{}
  {导师：吕建教授、黄宇教授；博士论文：《分布数据一致性技术研究》。}
\cventry{2005 年 9 月 -- 2009 年 6 月}{计算机科学与技术，学士}{南京大学}{计算机学院}{}{}

\section{工作经历}
\cventry{2025 年 8 月 -- 至今}{副教授}{湖南大学}{信息科学与工程学院}{}{}
\cventry{2020 年 9 月 -- 2025 年 7 月}{特任副研究员}{南京大学}{软件学院}{}{}
\cventry{2020 年 1 月 -- 2020 年 8 月}{特任副研究员}{南京大学}{计算机学院}{}{}
\cventry{2017 年 1 月 -- 2019 年 12 月}{助理研究员}{南京大学}{计算机学院}{}{}

\section{奖项与荣誉}
\cventry{2017 年度}{中国计算机学会}{CCF 优秀博士学位论文奖}{}{}{}
\cventry{2024 年}{全国大学生计算机系统能力大赛——编译系统设计赛}{优秀指导教师}{}{}{}
\cventry{2022 年}{南京大学，校级}{东吴证券奖教金}{}{}{}
\cventry{2019 年}{南京大学计算机学院，《问题求解》课程}{我心目中的好课程}{}{}{}

\clearpage
\section{科研项目}
% 项目名称、编号、角色与时间沿用原简历，网站未提供完整项目清单。
\cventry[0.4em]{2025 年 1 月 -- 2028 年 12 月}{国家自然科学基金，面上项目}
  {事务型数据库系统执行历史高效验证问题研究}{主持}{}{批准号：62472214。}
\cventry[0.4em]{2024 年 9 月 -- 2027 年 8 月}{江苏省基础研究计划自然科学基金，面上项目}
  {数据库系统事务一致性验证问题研究}{主持}{}{项目编号：BK20242014。}
\cventry[0.4em]{2021 年 12 月 -- 2022 年 12 月}{CCF-腾讯犀牛鸟基金，滚动支持}
  {分布式事务一致性协议设计与验证技术研究}{主持}{}{批准号：RAGR20200124。}
\cventry[0.4em]{2020 年 10 月 -- 2021 年 12 月}{CCF-腾讯犀牛鸟基金}
  {分布事务一致性模型与协议分析}{主持}{}{批准号：RAGR20200124。}
\cventry[0.4em]{2018 年 1 月 -- 2020 年 12 月}{国家自然科学基金，青年科学基金项目}
  {面向分布式系统的复制数据类型理论与技术研究}{主持}{}{批准号：61702253。}
\cventry[0.4em]{2021 年 1 月 -- 2022 年 1 月}{CCF-阿里巴巴创新研究计划（CCF-AIR 青年基金）}
  {基于 RDMA 和 NVM 加速的 ParallelRaft 协议设计与验证技术研究}{参与，项目主要成员}{}{项目编号：ATA5DHZ42 10001。}

\section{教学情况}
\cventry[0.4em]{2026 -- 2027 学年秋季}{湖南大学计算机学院}{程序设计}{}{}
  {\href{https://njusecourse.feishu.cn/wiki/VUVYwE7ZPiZDPqkWCDAc4R9Tn4d?from=from_copylink}{课程讲义}}
\cventry[0.4em]{2025 -- 2026 学年春季}{湖南大学计算机学院}{计算机数学}{}{}
  {\href{https://github.com/courses-at-hnu-by-hfwei/math4cs-lectures}{课程网站}\quad
   \href{https://space.bilibili.com/479141149/lists/7582381?type=season}{授课视频}}
\cventry[0.4em]{2021 -- 2024 年秋季}{南京大学软件学院}{C 语言程序设计基础}{}{}
  {\href{http://docs.cpl.icu}{课程网站}\quad
   \href{https://space.bilibili.com/479141149/channel/seriesdetail?sid=490580}{公开视频与教学资源}}
\cventry[0.4em]{2020 -- 2022 年秋季；2023 -- 2024 年春季}{南京大学软件学院}{编译原理}{}{}
  {\href{http://docs.compilers.cpl.icu}{课程网站}\quad
   \href{https://space.bilibili.com/479141149/channel/seriesdetail?sid=490581}{公开视频与教学资源}}
\cventry[0.4em]{2021 年春季}{南京大学软件学院}{离散数学}{}{}
  {\href{https://github.com/courses-at-nju-by-hfwei/discrete-math-lectures}{讲义}\quad
   \href{https://space.bilibili.com/479141149/channel/collectiondetail?sid=154382}{公开视频}}
\cventry[0.4em]{2017 年春季 -- 2020 年春季}{南京大学计算机学院，拔尖班课程}{问题求解}{}{}
  {\href{https://github.com/courses-at-nju-by-hfwei/problem-solving-class-lectures}{讲义与习题资源}}
\cventry[0.4em]{2011 -- 2019 年}{南京大学计算机学院}{算法设计与分析（助教）}{}{}
  {\href{https://tinyurl.com/algorithm-ta-tutorial}{习题课视频}}

\section{学生培养}
\cvitem{2025 -- 2026 年}{腾讯犀牛鸟精英人才计划：1 人。}
\cvitem{2024 -- 2025 年}{腾讯犀牛鸟精英人才计划：1 人。}
\cvitem{2024 年}{全国大学生计算机系统能力大赛——编译系统设计赛：指导 2 支队伍获得全国总决赛三等奖、优胜奖。}
\cvitem{2024 年}{腾讯犀牛鸟开源人才培养计划：3 人，其中任务奖金获得者 2 人、优秀学生 1 人（本科生）。}
\cvitem{2023 年}{指导校级优秀本科毕业论文：1 人。}

\clearpage
\section{论文与学术成果}
\textbf{粗体}标注本人姓名，$\ast$ 标注通讯作者。
2021 年及之后仅展示 CCF-A 类论文；2020 年及以前的论文保留，待刊论文单列。
% 使用显式引用列表筛选；全部论文数据仍保留在 publications.bib 中。
% CCF 会议目录仅计 Full/Regular paper，不计 Demo/Workshop 等形式。
% 分类来源：https://www.ccf.org.cn/Academic_Evaluation/By_category/2022-12-02/780280.shtml
% A 类目录：ARCH_DCP_SS（TPDS、ATC）、TCSE_SS_PDL（SOSP、ISSTA、ASE、OOPSLA）、DM_CS（VLDB、SIGMOD、ICDE）。
% 原列表顺序与全部 38 个引用键均保留；注释条目可按需恢复。
\nocite{DAT:PVLDB2027}
% \nocite{TAP:Festschrift2027} % 非 CCF-A，Festschrift/Workshop。
\nocite{VeriStrong:PVLDB2026}
\nocite{I3DP:SOSP2026}
% \nocite{InvWeaver:arXiv2026} % 预印本。
% \nocite{ElleDuplicateValues:SCCP2026} % Workshop。
% \nocite{SemanticConformance:SCCP2026} % Workshop。
% \nocite{Predicates:SCCP2026} % Workshop 扩展摘要。
\nocite{Clause2Inv:ISSTA2025}
\nocite{TimeKiller:ICDE2025}
\nocite{MiniTesting:ICDE2025}
\nocite{Specs:TPDS2025}
\nocite{LaM4Inv:ASE2024}
\nocite{Plume:OOPSLA2024}
% \nocite{IsoVista:PVLDB2024} % Demo Track，不在 CCF 推荐会议论文范围内。
\nocite{NOCNOC:SIGMOD2024}
% \nocite{MET:JSEP2024} % 非 CCF-A。
\nocite{PolySI:PVLDB2023}
% \nocite{TCC:ICPADS2022} % 非 CCF-A。
% \nocite{IncrementalCC:Internetware2022} % 非 CCF-A。
% \nocite{MongoDBCCC:JCST2022} % 非 CCF-A。
% \nocite{RemoveWin:ICPADS2021} % 非 CCF-A。
% \nocite{ByzGentleRain:SSS2021} % 非 CCF-A，Brief Announcement。
\nocite{UniStore:ATC2021}
% \nocite{PRaft:JOS2021} % 中文期刊，未列入国际 CCF-A 目录。
% \nocite{MongoDBCCC:Internetware2021} % 非 CCF-A。
\nocite{ASC:TPDS2021}
% 2020 年及以前不作 CCF-A 筛选。
\nocite{Jupiter:JCST2020}
\nocite{SpecFramework:SRDS2020}
\nocite{FastAtomic:PODC2020}
\nocite{TPaxos:JOS2020}
\nocite{CRDT-TLA:JOS2020}
\nocite{JupiterProtocol:OPODIS2018}
\nocite{JupiterProtocol:PODC2018}
\nocite{PA2AM:TC2017}
\nocite{RVSI:SRDS2017}
\nocite{VPC:TPDS2016}
\nocite{CTL3:PerCom2012}
\subsection{已发表的期刊与会议论文}
\printbibliography[heading=none,keyword=published]
\subsection{已接收、待刊论文}
\printbibliography[heading=none,keyword=forthcoming]
% \subsection{研讨会论文与扩展摘要}
% \printbibliography[heading=none,keyword=workshop]
% \subsection{预印本}
% \printbibliography[heading=none,keyword=preprint]

\clearpage
\section{学术服务}
\subsection{会议组织与程序委员会}
\cventry{2026 年}{首届一致性检查原理研讨会，与 VLDB 2026 同期举行}
  {\href{https://www.vldb.org/2026/workshops.html}{SCCP 2026 联合主席（Co-Chair）}}{}{}{}
\cventry{2026 年}{Session 4: Many Models, One Theory}
  {\href{https://sccp-workshop.github.io/sccp26/program.html}{SCCP 2026 分会场主席（Session Chair）}}{}{}{}
\cventry{2026 年}{第 45 届国际可靠分布式系统研讨会}
  {\href{https://srds-conference.org/index.php/program-committee/}{SRDS 2026 程序委员会委员}}{}{}{}
\cventry{2026 年}{第 13 届分布式数据一致性原理与实践研讨会}
  {\href{https://papoc-workshop.github.io/2026/committees.html}{PaPoC 2026 程序委员会委员}}{}{}{}

\subsection{研究成果评估与可复现性评审}
\cventry{2025 年}{Artifact Evaluation Committee}
  {\href{https://www.usenix.org/conference/atc25/call-for-artifacts}{USENIX ATC 2025 研究成果评估委员会委员}}{}{}{}
\cventry{2025 年}{Artifact Evaluation Committee}
  {\href{https://www.usenix.org/conference/osdi25/call-for-artifacts}{OSDI 2025 研究成果评估委员会委员}}{}{}{}
\cventry{2024、2025 年}{Availability \& Reproducibility Committee（ARC）}
  {\href{https://reproducibility.sigmod.org/2025/}{SIGMOD 可获取性与可复现性委员会委员}}{}{}{}

\subsection{期刊审稿}
\cvitem{TCS}{\href{https://www.sciencedirect.com/journal/theoretical-computer-science}{Theoretical Computer Science}}
\cvitem{JPDC}{\href{https://www.sciencedirect.com/journal/journal-of-parallel-and-distributed-computing}{Journal of Parallel and Distributed Computing}}
\cvitem{JSEP}{\href{https://onlinelibrary.wiley.com/journal/20477481}{Journal of Software: Evolution and Process}}
\cvitem{TAAS}{\href{https://dl.acm.org/journal/taas}{ACM Transactions on Autonomous and Adaptive Systems}}

\subsection{科普与人才培养报告}
\cvitem{2023 年}{“犀牛鸟中学科学人才培养计划”报告。}

\section{行政服务}
\cventry{2022 -- 2025 年}{南京大学软件学院}{教学委员会秘书}{}{}{}
\end{document}
